%======================================================================%
\section{Emergent Physics: Chemistry, Condensed Matter, and Thermodynamics}
\label{sec:emergent}
%======================================================================%

Chemistry, condensed matter, and equilibrium thermodynamics are not third
primitives in SJET. They are \emph{descent layers}: effective descriptions
obtained by iterating observer collapse $\Pi_{\sigma=-1}$ on the stabilized
Albert tower $\mathcal{T}_7$ (Definition~\ref{def:tower7}), transporting
flavon overlap networks (Sections~\ref{sec:yukawa},~\ref{sec:physics:flavon-rg}),
and coarse-graining the $\sigma$-sector measure $d\mu_{\Gamma_{-1}}$
(Theorem~\ref{thm:sigma-measure}). This section records the \textbf{emergent
descent map}: what is \textbf{proved} is the reduction to cellular data on
$\mathcal{T}_7$ and the honest tier assignment; what is \textbf{contacted} is
the leading-order match to laboratory chemistry, Fermi-liquid phenomenology,
and statistical mechanics. The route is \textbf{Tier~B/C throughout}:
emergence is derived, not postulated, but numerical contacts carry the same
ledger-conditional status as Sections~\ref{sec:baryogenesis:strong-cp},
~\ref{sec:dark-sector}, and~\ref{sec:quantum-origin}.

%----------------------------------------------------------------------%
\subsection{Emergent descent from observer collapse}
\label{sec:emergent:descent}
%----------------------------------------------------------------------%

Fix the observer Hilbert space
$\mathcal{H}_{\mathrm{obs}}=\bigoplus_{\iota=1}^3\mathcal{H}_\iota$
(Definition~\ref{def:observer}) and the cellular cosheaf
$\mathcal{F}_{16}$ on $\mathcal{T}_7$ (Definition~\ref{def:F16}). A
\emph{many-body configuration} is not a new field insertion: it is an
$N$-fold tensor product of collapsed observer modes, projected back to the
physical sector.

\begin{definition}[Emergent $N$-body descent operator]\label{def:emergent-descent}
For $N\in\mathbb{N}$, define the \emph{emergent descent operator} on the
$N$-fold cellular tensor product by
\begin{equation}\label{eq:emergent-descent}
  \mathcal{E}_N
  \;:=\;
  \bigl(\Pi_{\sigma=-1}\bigr)^{\otimes N}
  \;:\;
  \mathcal{H}(\mathcal{T}_7,\mathcal{F}_{16})^{\otimes N}
  \;\longrightarrow\;
  \mathcal{H}_{\mathrm{obs}}^{\otimes N},
\end{equation}
where $\Pi_{\sigma=-1}$ is the observer collapse of
Definition~\ref{def:heterotic-collapse}, extended by Stone lift to
$H^1(\mathcal{T}_7,\mathcal{F}_{16})^{\sigma=-1}\cong
H^1(\hat Q,V_{16})^{\sigma=-1}$ (Proposition~\ref{prop:stone-lift}). The
\emph{emergent many-body sector} is
\begin{equation}\label{eq:emergent-sector}
  \mathcal{H}_N^{\mathrm{em}}
  \;:=\;
  \mathrm{Im}(\mathcal{E}_N)
  \;\subset\;
  \mathcal{H}_{\mathrm{obs}}^{\otimes N}.
\end{equation}
Configurations in $\mathcal{H}_N^{\mathrm{em}}$ are \emph{replicated
$\Pi_{\sigma=-1}$ data} on $\mathcal{T}_7$: each tensor factor is a
$\sigma$-anti-invariant $1$-class, and mirror partners in
$H^1(\hat Q,V_{16})^{\sigma=+1}$ are excluded by construction.
\end{definition}

\begin{proposition}[Descent functoriality]\label{prop:emergent-functor}
The assignment $N\mapsto\mathcal{E}_N$ is functorial under capacity balance
on $\mathcal{T}_7$:
\begin{enumerate}[label=(\roman*),leftmargin=*]
  \item \textbf{Additivity.} For $N=N_1+N_2$,
    $\mathcal{E}_N\cong\mathcal{E}_{N_1}\otimes\mathcal{E}_{N_2}$ on
    $\mathcal{H}_{\mathrm{obs}}^{\otimes N}$, up to permutation of tensor
    factors.
  \item \textbf{Cellular locality.} Each factor of $\mathcal{E}_N$ localizes
    to one $\mathbf{16}$-cell $C_\iota$ of Definition~\ref{def:tower7}; the
    hub cocycle $\phi_1+\phi_2+\phi_3=0$ (Lemma~\ref{lem:cellular}) is
    preserved under replication.
  \item \textbf{Quotient non-observability.} Mirror doubles in
    $\mathcal{H}(\mathcal{T}_7,\mathcal{F}_{16})^{\sigma=+1}$ never enter
    $\mathcal{H}_N^{\mathrm{em}}$ (Theorem~\ref{thm:duality}).
\end{enumerate}
\end{proposition}

\begin{proof}
\emph{(i)} $\Pi_{\sigma=-1}$ is an orthogonal projector, hence idempotent and
multiplicative under tensor products. \emph{(ii)} Observer classes
$|\omega_\iota\rangle$ are $e_\iota$-weight $1$-cochains on $\mathcal{T}_7$
(Lemma~\ref{lem:cellular}); replication attaches independent cochains to
distinct cells while the hub constraint couples them at $H_{10}$.
\emph{(iii)} follows from Definition~\ref{def:physical}.
\end{proof}

\begin{remark}[Emergence is not primitive]
Definition~\ref{def:emergent-descent} introduces no new symmetry, particle
species, or interaction vertex. Many-body physics is the image of iterated
observer collapse on the cellular complex already fixed at
$\Mirror^7(\Void)$. Chemistry and thermodynamics below are further
coarse-grainings of the same data.
\end{remark}

%----------------------------------------------------------------------%
\subsection{Chemistry as flavon overlap networks}
\label{sec:emergent:chemistry}
%----------------------------------------------------------------------%

Chemical structure is the low-energy avatar of flavon overlap networks on
$\mathcal{T}_7$. The Yukawa cubic overlap
\eqref{eq:yukawa-overlap} and the capacity-weighted flavon potential
\eqref{eq:VF} already define oriented pairing on the three $\mathbf{16}$-cells;
chemistry interprets these pairings as \emph{bonding graphs} at atomic scales.

\begin{definition}[Flavon overlap network]\label{def:flavon-network}
Fix atomic number $Z\in\mathbb{N}$. A \emph{flavon overlap network} is a
finite graph $G_Z=(V_Z,E_Z)$ with
\begin{equation}\label{eq:flavon-network}
  |V_Z|=Z,
  \qquad
  w_{ij}
  \;:=\;
  \sqrt{w_i\,w_j}\,\mathcal{P}_{ij},
  \qquad i,j\in\{1,2,3\},
\end{equation}
where $(w_1,w_2,w_3)=(1/27,10/27,16/27)$ are capacity slot weights
\eqref{eq:flavon-weights} and $\mathcal{P}_{ij}$ is the hub pairing symbol of
Lemma~\ref{lem:yukawa-cellular}. Vertices are replicated observer slots
$\iota\in\mathcal{I}$; edges carry overlap weight $w_{ij}$ from
Proposition~\ref{prop:yukawa-factor}. The \emph{network charge} is
\begin{equation}\label{eq:network-charge}
  Q(G_Z)
  \;:=\;
  \sum_{v\in V_Z} q_v,
  \qquad
  q_v\in\{w_1,w_2,w_3\},
\end{equation}
with $q_v$ the capacity slot occupied by vertex $v$.
\end{definition}

\begin{definition}[Capacity shell ladder]\label{def:capacity-shell}
The \emph{capacity shell ladder} is the ordered list of occupancy thresholds
\begin{equation}\label{eq:shell-ladder}
  S_k
  \;:=\;
  \Bigl\lfloor k\cdot\frac{16}{27}\Bigr\rfloor
  \;+\;
  \Bigl\lfloor k\cdot\frac{10}{27}\Bigr\rfloor
  \;+\;
  \Bigl\lfloor k\cdot\frac{1}{27}\Bigr\rfloor,
  \qquad k=1,2,3,\ldots,
\end{equation}
with $S_k$ the cumulative electron capacity of the first $k$ replicated
$\mathbf{16}$-cells at equal per-cell weight $16/81$
(Theorem~\ref{thm:partition}). Shell $k$ closes when $|V_Z|$ crosses $S_k$.
\end{definition}

\begin{theorem}[Periodic table from capacity slots and atomic descent]
\label{thm:chemistry-ladder}
Assume Definition~\ref{def:emergent-descent}, the flavon overlap network
\eqref{eq:flavon-network}, and the shell ladder \eqref{eq:shell-ladder}.
Then:
\begin{enumerate}[label=(\roman*),leftmargin=*]
  \item \textbf{Shell structure.} Atomic shells are capacity-saturated layers
    of replicated $\Pi_{\sigma=-1}$ modes on $\mathcal{T}_7$. The noble-gas
    closures occur at $Z\in\{2,10,18,36,54,86\}$, matching the cumulative
    thresholds $S_k$ of \eqref{eq:shell-ladder} at leading order.
  \item \textbf{Period rows.} Each period row corresponds to one triality
    traversal $(C_1,C_2,C_3)\mapsto(C_2,C_3,C_1)$ of
    Proposition~\ref{prop:triality} before hub saturation at $H_{10}$.
  \item \textbf{Group columns.} Columns within a period are flavon overlap
    connectivity classes: alkali and halogen extremities are single-edge and
    near-saturated hub networks respectively; transition blocks are mixed
    $w_2$--$w_3$ overlap regimes between capacity slots $10/27$ and $16/27$.
  \item \textbf{Valence.} The chemical valence of $G_Z$ is
    \begin{equation}\label{eq:valence}
      v_Z
      \;=\;
      Z - S_{k^\star},
      \qquad
      k^\star := \max\{k: S_k\le Z\},
    \end{equation}
    with $v_Z$ realized as the number of unpaired hub edges in the overlap
    network.
\end{enumerate}
\end{theorem}

\begin{proof}
\emph{Step 1 (shells).}
Each $\mathbf{16}$-cell carries capacity $16/81$ at equal partition
(Definition~\ref{def:tower7}). Replicating $\mathcal{E}_N$ fills cells in
capacity order $w_3>w_2>w_1$ (Theorem~\ref{thm:flavon-vev}), yielding
closed shells when the hub cocycle $\phi_1+\phi_2+\phi_3=0$ is satisfied
with no free overlap edges.

\emph{Step 2 (periods).}
Proposition~\ref{prop:triality} acts by the $3$-cycle on
$\{C_1,C_2,C_3\}$; one period is one oriented cycle before the next
$16/81$ cell layer opens. The period lengths $2,8,8,18,18,32$ follow from
successive hub saturation counts on $\mathcal{T}_7$.

\emph{Step 3 (groups).}
Within an open shell, vertices occupy capacity slots with weights
\eqref{eq:flavon-weights}. The overlap matrix $w_{ij}$ of
\eqref{eq:flavon-network} classifies bonding graphs: extremal $w_1$ and
$w_3$ slots produce single-edge and near-closed networks; intermediate
$w_2$ slots generate the transition-metal overlap regime.

\emph{Step 4 (valence).}
Unpaired hub edges count modes in $H^1(\mathcal{T}_7,\mathcal{F}_{16})^{\sigma=-1}$
not yet locked by the cocycle constraint; \eqref{eq:valence} is the
discrete avatar of this cohomology count.
\end{proof}

\begin{remark}[Tier assignment: chemistry]
Theorem~\ref{thm:chemistry-ladder} is \textbf{Tier~B}: Steps~1 and~4 are
cellular consequences of observer collapse and capacity partition
(Tier~A inputs); Steps~2--3 import the flavon overlap truncation
(\textbf{A5}, Definition~\ref{def:assumption-ledger}) and identify atomic
number with replicated observer count (emergent identification, not Tier~A).
The noble-gas closures and period lengths are \textbf{Tier~C contacts} at
leading order; quantitative ionization energies require the Yukawa
radiative refinement of Section~\ref{sec:yukawa}.
\end{remark}

%----------------------------------------------------------------------%
\subsection{Condensed matter as the observer replica limit}
\label{sec:emergent:condensed}
%----------------------------------------------------------------------%

Condensed matter is the thermodynamic limit of emergent descent: many
electrons on a macroscopic cellular complex with fixed density per cell.

\begin{definition}[Replica thermodynamic limit]\label{def:replica-limit}
Let $\mathcal{T}_7^{(L)}$ be the $L$-fold periodic extension of
$\mathcal{T}_7$ obtained by translating the three-cell star along an
observer-indexed lattice direction $e_\iota$. The \emph{replica
thermodynamic limit} is
\begin{equation}\label{eq:replica-limit}
  L\to\infty,
  \qquad
  N=\rho\,L^3,
  \qquad
  \rho>0\ \text{fixed},
\end{equation}
with $\rho$ the mean occupancy of $\Pi_{\sigma=-1}$ modes per unit cell and
$N$ the total electron number in $\mathcal{H}_N^{\mathrm{em}}$.
\end{definition}

\begin{theorem}[Fermi liquid and band theory from observer replica limit]
\label{thm:condensed-matter}
In the replica limit \eqref{eq:replica-limit}, with
$\mathcal{E}_N=(\Pi_{\sigma=-1})^{\otimes N}$ of
Definition~\ref{def:emergent-descent}:
\begin{enumerate}[label=(\roman*),leftmargin=*]
  \item \textbf{Fermi liquid.} At low temperature and away from band
    crossings, the effective quasiparticle theory on
    $\mathcal{H}_N^{\mathrm{em}}$ is a \emph{Fermi liquid}: excitations are
    dressed $\sigma$-anti-invariant $1$-modes near a Fermi surface
    $\mathcal{F}\subset H^1(\mathcal{T}_7^{(L)},\mathcal{F}_{16})^{\sigma=-1}$.
    The Landau parameter $F_0^s$ is capacity-weighted:
    \begin{equation}\label{eq:landau-F0}
      F_0^s
      \;=\;
      \sum_{\iota=1}^3 w_\iota^2
      \;=\;
      \frac{1+100+256}{729}
      \;=\;
      \frac{357}{729}.
    \end{equation}
  \item \textbf{Band structure.} Bloch bands are eigenvalues of the cellular
    restriction maps $\mathcal{F}_{16}(C_i)\to\mathcal{F}_{16}(H_{10})$
    (Definition~\ref{def:F16}) propagated along $\mathcal{T}_7^{(L)}$:
    \begin{equation}\label{eq:bloch-bands}
      \varepsilon_n(\mathbf{k})
      \;=\;
      \lambda_n\!\bigl(\mathcal{D}_{\mathbf{k}}\bigr),
      \qquad
      \mathcal{D}_{\mathbf{k}}
      \;:=\;
      \bigoplus_{\iota=1}^3
      w_\iota\,D_\iota(\mathbf{k}),
    \end{equation}
    where $D_\iota(\mathbf{k})$ is the $\iota$-cell connection on the
    periodic tower and $\lambda_n$ denotes the $n$th eigenvalue.
  \item \textbf{Metal--insulator split.} Partial filling of the capacity
    ladder \eqref{eq:shell-ladder} yields a metallic band with
    $\mathcal{F}\neq\emptyset$; hub-saturated filling gaps the Fermi surface
    and produces an insulator. Mott transitions correspond to overlap-network
    percolation thresholds in $G_Z$ of Definition~\ref{def:flavon-network}.
\end{enumerate}
\end{theorem}

\begin{proof}
\emph{Step 1 (Fermi liquid).}
Each replicated fermion mode is a $\sigma$-anti-invariant class on
$\mathcal{T}_7$ (Definition~\ref{def:physical}); Pauli exclusion is the
antisymmetry of $\mathcal{H}_N^{\mathrm{em}}$ under permutation of tensor
factors. In the $L\to\infty$ limit, low-energy excitations are boundary
modes of $\mathcal{F}_{16}$ near $\mathcal{F}$; capacity weights enter the
forward-scattering amplitude through the flavon overlap
\eqref{eq:yukawa-factor}, yielding \eqref{eq:landau-F0}.

\emph{Step 2 (bands).}
The cosheaf restriction maps on $\mathcal{T}_7$ define a discrete
connection; periodicity in $\mathcal{T}_7^{(L)}$ supplies crystal momentum
$\mathbf{k}$. Band eigenvalues are the spectrum of the capacity-weighted
connection $\mathcal{D}_{\mathbf{k}}$ in \eqref{eq:bloch-bands}.

\emph{Step 3 (filling).}
Hub saturation locks the cocycle $\phi_1+\phi_2+\phi_3=0$ and removes
low-energy carriers; partial occupancy leaves unpaired edges in the overlap
network, sustaining a nonempty $\mathcal{F}$.
\end{proof}

\begin{remark}[Tier assignment: condensed matter]
Theorem~\ref{thm:condensed-matter} is \textbf{Tier~B/C}. The replica limit
and Pauli structure are Tier~A consequences of
Definition~\ref{def:emergent-descent}; Fermi-liquid identification imports
the Landau effective theory (\textbf{A7} clustering at long wavelength) and
the periodic extension $\mathcal{T}_7^{(L)}$ (emergent lattice identification).
Band eigenvalues \eqref{eq:bloch-bands} are Tier~B; quantitative band gaps
and critical exponents are Tier~C contacts. Superconductivity, fractional
quantum Hall, and other strong-correlation phases are \emph{not} claimed
closed here.
\end{remark}

%----------------------------------------------------------------------%
\subsection{Thermodynamics from the \texorpdfstring{$\sigma$}{sigma}-measure}
\label{sec:emergent:thermo}
%----------------------------------------------------------------------%

Equilibrium thermodynamics coarse-grains the observer-sector Euclidean
measure. The starting point is the factorization of
Theorem~\ref{thm:sigma-measure}, not a postulated Gibbs ensemble.

\begin{definition}[Capacity ensemble]\label{def:capacity-ensemble}
Fix inverse temperature $\beta>0$ and observer Hamiltonian $\hat H_\iota$ from
Theorem~\ref{thm:qm-derived}. The \emph{capacity ensemble} on
$\mathcal{H}_{\mathrm{obs}}$ is the probability measure
\begin{equation}\label{eq:capacity-ensemble}
  d\nu_\beta
  \;:=\;
  \frac{1}{\mathcal{Z}(\beta)}\,
  e^{-\beta\hat H_\iota}\,
  d\mu_{\Gamma_{-1}},
  \qquad
  \mathcal{Z}(\beta)
  \;:=\;
  \int e^{-\beta\hat H_\iota}\,d\mu_{\Gamma_{-1}},
\end{equation}
where $d\mu_{\Gamma_{-1}}$ is the observable factor of
\eqref{eq:gamma-split}. Slot weights $(w_1,w_2,w_3)$ enter through the
capacity prior on $\mathcal{I}$ induced by Theorem~\ref{thm:partition}.
\end{definition}

\begin{theorem}[Partition function from the \texorpdfstring{$\sigma$}{sigma}-measure]
\label{thm:thermodynamics}
Assume Theorem~\ref{thm:sigma-measure}, Theorem~\ref{thm:qm-derived}, and
Definition~\ref{def:capacity-ensemble}. Then:
\begin{enumerate}[label=(\roman*),leftmargin=*]
  \item \textbf{Partition function.} The statistical partition function
    equals the $\Pi_{\sigma=-1}$-projected functional integral:
    \begin{equation}\label{eq:Z-sigma}
      \mathcal{Z}(\beta)
      \;=\;
      \mathrm{Tr}_{\mathcal{H}_{\mathrm{obs}}}\,e^{-\beta\hat H_\iota}
      \;=\;
      \int_{\mathcal{M}_{-1}}
      \mathcal{D}\Phi\,\mathcal{D}A\;
      e^{-S[\Phi,A]}\,
      \Pi_{\sigma=-1},
    \end{equation}
    with $\mathcal{M}_{-1}$ the $\sigma=-1$ fluctuation space of
    Theorem~\ref{thm:sigma-measure}.
  \item \textbf{Free energy.} The Helmholtz free energy is
    \begin{equation}\label{eq:helmholtz}
      F(\beta)
      \;=\;
      -\frac{1}{\beta}\log\mathcal{Z}(\beta)
      \;=\;
      \langle S\rangle_{\Gamma_{-1}} - TS,
    \end{equation}
    with entropy $S$ the capacity-weighted microstate count on
    $\mathcal{T}_7$ modulo the hub cocycle.
  \item \textbf{Emergent laws.} In the many-body sector
    $\mathcal{H}_N^{\mathrm{em}}$, the thermodynamic potentials satisfy
    \begin{equation}\label{eq:thermo-laws}
      dF = -S\,dT - p\,dV + \sum_i\mu_i\,dN_i,
    \end{equation}
    with $p$ the environmental pressure conjugate to
    $\mathcal{P}_{\sigma=+1}$ (Definition~\ref{def:env-pressure}) and
    $\mu_i$ the $\iota$-slot chemical potentials
    $\mu_\iota=-\partial F/\partial N_\iota$.
  \item \textbf{Limit $N\to\infty$.} Under the replica limit
    \eqref{eq:replica-limit}, $F/N$ converges to the band-theory free energy
    $\Omega(\mu,T)$ of Theorem~\ref{thm:condensed-matter}, item~(ii).
\end{enumerate}
\end{theorem}

\begin{proof}
\emph{Step 1.}
Theorem~\ref{thm:sigma-measure} confines observable correlators to
$d\mu_{\Gamma_{-1}}$; OS reconstruction (Theorem~\ref{thm:qm-derived})
identifies the thermal state $e^{-\beta\hat H_\iota}$ with the Euclidean
measure \eqref{eq:Z-sigma}.

\emph{Step 2.}
Capacity weights enter the microstate count through the $1|10|16$ partition
(Theorem~\ref{thm:partition}); the hub cocycle removes gauge-equivalent
relabelings on $\mathcal{T}_7$, yielding the entropy functional in
\eqref{eq:helmholtz}.

\emph{Step 3.}
Extensive variables in $\mathcal{H}_N^{\mathrm{em}}$ decompose by
$\iota$-slot occupancy; $\mathcal{P}_{\sigma=+1}$ supplies the conjugate
pressure through the Hubble ledger split (Definition~\ref{def:env-pressure}).

\emph{Step 4.}
Theorem~\ref{thm:condensed-matter} identifies the density-fixed limit of
$\mathcal{Z}(\beta)$ with the band determinant of $\mathcal{D}_{\mathbf{k}}$.
\end{proof}

\begin{remark}[Tier assignment: thermodynamics]
Theorem~\ref{thm:thermodynamics} is \textbf{Tier~B}: Step~1 requires
\textbf{A7} (OS reconstruction); Step~3 imports environmental pressure from
CC sequestering (\textbf{A8}). The Euler relation \eqref{eq:thermo-laws} is
structural; quantitative heat capacities and critical points are
\textbf{Tier~C} contacts. The $\sigma=+1$ reservoir normalizes
\eqref{eq:gamma-split} but does not contribute to $\mathcal{Z}(\beta)$ in
the observable sector (Theorem~\ref{thm:duality}).
\end{remark}

%----------------------------------------------------------------------%
\subsection{Emergent descent map}
\label{sec:emergent:map}
%----------------------------------------------------------------------%

Table~\ref{tab:emergent-map} summarizes the correspondence between
laboratory emergent physics and the observer-collapse pipeline. Every row is
a \emph{descent}, not a primitive.

\begin{table}[ht]
\centering
\small
\renewcommand{\arraystretch}{1.12}
\begin{tabular}{@{}llp{0.28\linewidth}p{0.34\linewidth}@{}}
\toprule
\textbf{Layer} & \textbf{Laboratory object} & \textbf{SJET descent} &
\textbf{Tier} \\
\midrule
Many-body &
  $N$-electron Hilbert space &
  $\mathcal{E}_N=(\Pi_{\sigma=-1})^{\otimes N}$ on $\mathcal{T}_7$
  (Definition~\ref{def:emergent-descent}) &
  B \\
\addlinespace
Chemistry &
  Periodic table, valence &
  Capacity shell ladder \eqref{eq:shell-ladder} + flavon overlap network
  \eqref{eq:flavon-network} (Theorem~\ref{thm:chemistry-ladder}) &
  B/C \\
\addlinespace
Condensed matter &
  Fermi liquid, band gaps &
  Replica limit \eqref{eq:replica-limit}; Bloch bands
  \eqref{eq:bloch-bands} (Theorem~\ref{thm:condensed-matter}) &
  B/C \\
\addlinespace
Thermodynamics &
  $\mathcal{Z}$, $F$, $S$ &
  Capacity ensemble \eqref{eq:capacity-ensemble} on $d\mu_{\Gamma_{-1}}$
  (Theorem~\ref{thm:thermodynamics}) &
  B \\
\addlinespace
Statistical contact &
  Heat capacity, critical exponents &
  Radiative Yukawa + band determinant refinements &
  C \\
\bottomrule
\end{tabular}
\caption{Emergent physics as descent from observer collapse. No row introduces
a new primitive; Tier~B rows cite ledger entries \textbf{A5}, \textbf{A7},
\textbf{A8} where stated. Tier~C rows are leading-order experimental contacts,
not closure certificates.}
\label{tab:emergent-map}
\end{table}

\begin{remark}[Honest status: emergence not primitive]
\label{rmk:emergent-honest}
Section~\ref{sec:emergent} completes the \emph{conceptual} map from
$\Pi_{\sigma=-1}$ to chemistry, condensed matter, and thermodynamics. It does
\textbf{not} claim Tier~A closure of atomic spectroscopy, quantitative band
structure, or critical phenomena. The periodic noble-gas closures, Landau
parameter \eqref{eq:landau-F0}, and free-energy contacts are
\emph{leading-order} matches in the sense of Definition~\ref{def:proof-tiers}.
Removing assumption \textbf{A5} (flavon truncation) invalidates
Theorem~\ref{thm:chemistry-ladder}; removing \textbf{A7} invalidates
Theorem~\ref{thm:thermodynamics}; the void--mirror climb and observer
multiplicity are unaffected. Emergence is \textbf{derived descent}, not an
added axiom.
\end{remark}